\documentclass[a4paper,11pt]{ctexart}
\usepackage{amsmath, amssymb, amsfonts}
\usepackage{geometry}
\geometry{left=1.5cm, right=1.5cm, top=2.0cm, bottom=2.0cm, headsep=0.3cm, footskip=1cm}
\usepackage{listings}
\usepackage{xcolor}
\usepackage{tikz}
\usepackage{pgfplots}
\usepackage{hyperref}
\usepackage{fancyhdr}
\usepackage{tcolorbox}
\tcbuselibrary{breakable, skins}
\usepackage{array}
\usepackage{booktabs}
\usepackage[shortlabels]{enumitem}
\usepackage{needspace}
\usepackage{multirow}

\AtBeginDocument{
  \hypersetup{colorlinks=true, linkcolor=blue, filecolor=magenta, urlcolor=cyan}
}

\setlist{nosep, itemsep=0.8ex, parsep=0.1em}
\setlist[itemize]{leftmargin=1.8em}
\setlist[enumerate]{leftmargin=1.8em}

% ===== 颜色定义 =====
\definecolor{codegreen}{rgb}{0,0.6,0}
\definecolor{codegray}{rgb}{0.5,0.5,0.5}
\definecolor{codepurple}{rgb}{0.58,0,0.82}
\definecolor{codebg}{rgb}{0.98,0.98,0.98}
\definecolor{tipsblue}{rgb}{0.85,0.93,1.0}
\definecolor{highlightyellow}{rgb}{1.0,0.97,0.8}

% ===== 代码块样式 =====
\lstdefinestyle{mystyle}{
    backgroundcolor=\color{codebg},
    commentstyle=\color{codegreen},
    keywordstyle=\color{magenta}\bfseries,
    numberstyle=\tiny\color{codegray},
    stringstyle=\color{codepurple},
    basicstyle=\ttfamily\small,
    breakatwhitespace=false,
    breaklines=true,
    captionpos=b,
    keepspaces=true,
    numbers=left,
    numbersep=6pt,
    showspaces=false,
    showstringspaces=false,
    showtabs=false,
    tabsize=2,
    language=Python,
    frame=single,
    framerule=0.5pt,
    rulecolor=\color{gray!40}
}
\lstset{style=mystyle}

% ===== 彩色框 =====
\newtcolorbox{problembox}[1][]{
    colback=white,
    colframe=blue!60!black,
    title=\textbf{#1},
    fonttitle=\bfseries\small,
    boxrule=1pt,
    arc=3pt, outer arc=3pt,
    coltitle=black,
    before skip=10pt, after skip=8pt
}

\newtcolorbox{solutionbox}[1][]{
    enhanced, breakable,
    colback=green!3!white,
    colframe=green!50!black,
    title=\textbf{#1},
    fonttitle=\bfseries\small,
    boxrule=1pt,
    arc=3pt, outer arc=3pt,
    coltitle=black,
    before skip=8pt, after skip=10pt
}

\newtcolorbox{metaphorbox}[1][]{
    colback=orange!5!white,
    colframe=orange!70!black,
    title=\textbf{#1},
    fonttitle=\bfseries\small,
    boxrule=1pt,
    arc=3pt, outer arc=3pt,
    coltitle=black,
    before skip=8pt, after skip=8pt
}

\newtcolorbox{beginnerbox}[1][]{
    colback=tipsblue,
    colframe=blue!50!black,
    title=\textbf{#1},
    fonttitle=\bfseries\small,
    boxrule=1pt,
    arc=3pt, outer arc=3pt,
    coltitle=black,
    before skip=8pt, after skip=8pt
}

\newtcolorbox{appbox}[1][]{
    colback=green!3!white,
    colframe=green!60!black,
    title=\textbf{#1},
    fonttitle=\bfseries\small,
    boxrule=1pt,
    arc=3pt, outer arc=3pt,
    coltitle=black,
    before skip=8pt, after skip=10pt
}

\newtcolorbox{keybox}[1][]{
    colback=highlightyellow,
    colframe=orange!80!black,
    title=\textbf{#1},
    fonttitle=\bfseries\small,
    boxrule=1pt,
    arc=3pt, outer arc=3pt,
    coltitle=black,
    before skip=8pt, after skip=8pt
}

% ===== 页眉页脚 =====
\pagestyle{fancy}
\fancyhf{}
\fancyhead[R]{深度学习-第3章}
\fancyhead[L]{神经网络的学习 · 练习题}
\fancyfoot[C]{\thepage}
\addtolength{\headheight}{2pt}

\parskip=2pt plus 1pt minus 0.5pt
\linespread{1.35}
\raggedbottom
\widowpenalty=10000
\clubpenalty=10000

\title{\textbf{深度学习\\第3章\quad 神经网络的学习}\\[0.5em]
{\large\color{gray}——Loss、梯度、SGD：搞懂这三件事，神经网络就真的会"学"了}}
\date{\today}

\begin{document}

\maketitle

% ===== 章节导读 =====
\begin{beginnerbox}[本章题目导览：你将挑战哪些核心问题？]
    前两章我们搭出了神经网络的"骨架"——知道数据怎么流过网络、激活函数是什么。但网络还不会"学"，它不知道自己预测得好不好，也不知道怎么变得更好。本章解决这两个问题，共 6 道题：

    \begin{itemize}
        \item \textbf{Q3-01（辨析填空）}：常见损失函数横向对比——MSE、MAE、交叉熵、Huber Loss，各自的公式、适用任务和主要特点是什么？
        \item \textbf{Q3-02（计算题）}：损失函数手动计算——给定预测值和真实标签，分别计算 MSE 损失和交叉熵损失，把公式变成具体数字，感受"哪种惩罚更重"。
        \item \textbf{Q3-03（概念理解）}：导数、偏导数与梯度——三个概念的联系与区别，以及数值微分的原理和计算。
        \item \textbf{Q3-04（推导计算）}：神经网络的梯度计算——对一个简单的单层网络，手动用数值微分计算损失对权重的梯度，理解梯度的物理含义。
        \item \textbf{Q3-05（对比分析）}：梯度下降三兄弟——批量梯度下降（BGD）、随机梯度下降（SGD）、小批量梯度下降（Mini-batch GD）的区别、优缺点和适用场景。
        \item \textbf{Q3-06（综合代码实现）}：训练循环全流程——阅读并补全一个含前向传播、损失计算、梯度更新的完整训练代码，追踪每一步的数据变化。
    \end{itemize}
\end{beginnerbox}

\section{第 3 章\quad 神经网络的学习}

前两章搭好了网络的骨架，但它还不会"学"。给它一张猫的图片，它输出一个随机的预测值，反正也不知道对不对，更不知道怎么改进。

要让网络"学起来"，需要解决三件事：

\begin{itemize}
    \item \textbf{怎么衡量"预测有多差"？}——损失函数（Loss Function）
    \item \textbf{怎么知道"往哪个方向调整参数"？}——梯度（Gradient）
    \item \textbf{怎么一步步调整让 Loss 变小？}——梯度下降（Gradient Descent）
\end{itemize}

这三件事合起来，就是神经网络"学习"的完整闭环。本章把这个闭环从数学到代码全部打通。

\begin{metaphorbox}[先建立一个总感觉：学习是什么过程？]
    \textbf{把神经网络的学习比作一个蒙眼射箭手在雾中找靶心。}

    \begin{itemize}
        \item \textbf{损失函数} = 报靶员：射完一箭后，告诉你"偏了多少米"——这就是 Loss，Loss 越小越好；
        \item \textbf{梯度} = 指南针：告诉你"要让 Loss 减小，弓应该往哪个方向转"——梯度指向 Loss 增大最快的方向，所以朝梯度的反方向调整；
        \item \textbf{梯度下降} = 每次微调弓的角度：根据指南针，往减小 Loss 的方向小步移动，一轮一轮地迭代，越来越接近靶心；
        \item \textbf{SGD} = 每次不用把所有箭都射一遍再调整，而是射一箭（或一小批箭）就调整一次——虽然每次估计可能有误差，但调整频率更高，总体反而更快。
    \end{itemize}
\end{metaphorbox}


%% ============================================================
\Needspace{5\baselineskip}
\begin{problembox}[Q3-01\quad 常见损失函数横向对比]
    \textbf{题目描述：}\\
    文档第 3.1 节介绍了多种损失函数。损失函数是"评分系统"——选对了，模型才能朝正确方向优化；选错了，优化方向就会偏离目标。

    \vspace{0.3cm}
    \textbf{问题：}
    \begin{enumerate}[(1)]
        \item 填写下表，梳理四种常见损失函数的关键信息：

        \begin{center}
        \begin{tabular}{p{2.2cm} p{4.5cm} p{2.5cm} p{2.5cm}}
            \toprule
            损失函数 & 数学公式 & 适用任务 & 主要特点 \\
            \midrule
            均方误差（MSE） & $\mathcal{L} = \dfrac{1}{N}\displaystyle\sum_{i=1}^N (\hat{y}_i - y_i)^2$ & \underline{\hspace{2cm}} & \underline{\hspace{2cm}} \\[2.0em]
            平均绝对误差（MAE） & $\mathcal{L} = \dfrac{1}{N}\displaystyle\sum_{i=1}^N |\hat{y}_i - y_i|$ & \underline{\hspace{2cm}} & \underline{\hspace{2cm}} \\[2.0em]
            交叉熵（CE） & $\mathcal{L} = -\dfrac{1}{N}\displaystyle\sum_{i=1}^N\sum_{k=1}^K t_{ik}\log \hat{y}_{ik}$ & \underline{\hspace{2cm}} & \underline{\hspace{2cm}} \\[2.0em]
            Huber Loss & $\mathcal{L}_\delta = \begin{cases}\frac{1}{2}(\hat{y}-y)^2 & |\hat{y}-y|\leq\delta \\ \delta|\hat{y}-y|-\frac{\delta^2}{2} & \text{otherwise}\end{cases}$ & \underline{\hspace{2cm}} & \underline{\hspace{2cm}} \\
            \bottomrule
        \end{tabular}
        \end{center}

        \item MSE 和 MAE 都用于回归任务，但对\textbf{离群点（outliers）}的敏感程度不同。解释为什么 MSE 比 MAE 对离群点更敏感，以及各自适合什么数据分布。

        \item 为什么分类任务通常使用交叉熵损失而不是 MSE？从梯度角度解释：当预测结果很差时（预测概率接近0而真实标签为1），两种损失函数的梯度有何差异？

        \item Huber Loss 被称为"MSE 和 MAE 的结合体"。解释为什么，以及超参数 $\delta$ 如何控制这种平衡。
    \end{enumerate}

    \vspace{0.2cm}
    {\color{gray}\small\textit{来源溯源：[文档第 3 章 3.1 损失函数]}}
\end{problembox}

\begin{beginnerbox}[小白必读：损失函数为什么不能随便选？]
    \begin{itemize}
        \item \textbf{损失函数决定了优化目标}：梯度下降会让 Loss 减小——你选什么 Loss，模型就朝那个方向优化。选 MSE，模型会努力让平均平方误差最小；选交叉熵，模型会努力让预测概率分布和真实分布接近。目标不同，训练出来的模型行为就不同。
        \item \textbf{分类为什么不用 MSE}：分类标签是 $\{0, 1\}$ 或 one-hot 向量，MSE 在两端（预测值接近0或1）梯度极小，训练极慢；而交叉熵在预测错误时（预测概率接近0但真实为1）梯度很大，能快速纠错。
        \item \textbf{损失函数要可微}：梯度下降依赖梯度，损失函数在大多数点必须可微。MAE 在零点不可微（但在实践中可用次梯度或 Huber 替代），准确率（Accuracy）完全不可微，因此不能直接作为损失函数，只能作为评估指标。
    \end{itemize}
\end{beginnerbox}

\begin{solutionbox}[参考答案与详细解析]
    \textbf{(1) 四种损失函数对比表：}

    \begin{center}
    \begin{tabular}{p{2.2cm} p{4.5cm} p{2.2cm} p{2.8cm}}
        \toprule
        损失函数 & 公式 & 适用任务 & 主要特点 \\
        \midrule
        MSE & $\frac{1}{N}\sum(\hat{y}_i-y_i)^2$ & 回归 & 平滑可微，对离群点敏感 \\[0.8em]
        MAE & $\frac{1}{N}\sum|\hat{y}_i-y_i|$ & 回归 & 对离群点鲁棒，零点不可微 \\[0.8em]
        交叉熵 & $-\frac{1}{N}\sum_i\sum_k t_{ik}\log\hat{y}_{ik}$ & 分类 & 配合 Softmax，梯度信号强 \\[0.8em]
        Huber & 小误差用 MSE，大误差用 MAE & 回归 & 两者折中，$\delta$ 控制边界 \\
        \bottomrule
    \end{tabular}
    \end{center}

    \textbf{(2) MSE vs MAE 对离群点的敏感度：}

    MSE 对误差\textbf{平方}，一个误差为10的离群点贡献 $10^2=100$，而误差为1的正常点贡献1。离群点对 MSE 的影响是正常点的100倍。MAE 对误差取绝对值，离群点贡献10，正常点贡献1，差距仅10倍。

    \begin{itemize}
        \item \textbf{数据无离群点/误差服从正态分布}：选 MSE（最大似然估计在正态噪声假设下等价于 MSE）；
        \item \textbf{数据含离群点/误差服从重尾分布}：选 MAE 或 Huber Loss。
    \end{itemize}

    \textbf{(3) 交叉熵 vs MSE 用于分类——梯度角度：}

    设真实标签 $y=1$，预测概率 $\hat{y} \approx 0$（预测严重错误）。

    \textbf{交叉熵梯度}：$\partial \mathcal{L}/\partial \hat{y} = -1/\hat{y}$，当 $\hat{y} \to 0$ 时梯度趋向无穷大——\textbf{预测越差，纠错信号越强}，收敛快。

    \textbf{MSE 梯度}：$\partial \mathcal{L}/\partial \hat{y} = 2(\hat{y}-y) \approx -2$，无论预测多差，梯度幅度几乎不变，纠错信号微弱，训练极慢。此外，Sigmoid/Softmax 输出层与 MSE 组合后，梯度中会出现 $\sigma'(z)$ 项（最大值仅0.25），进一步削弱梯度。

    \textbf{(4) Huber Loss 的双重特性：}

    \begin{itemize}
        \item 误差较小时（$|\hat{y}-y| \leq \delta$）：退化为 MSE（$\frac{1}{2}e^2$），平滑可微，梯度稳定；
        \item 误差较大时（超过 $\delta$）：退化为 MAE（线性惩罚），不再随误差平方增长，抑制离群点影响。
    \end{itemize}

    $\delta$ 控制"什么程度算大误差"：$\delta$ 越大，越接近纯 MSE；$\delta$ 越小，越接近纯 MAE。实践中 $\delta$ 通常通过交叉验证选取。
\end{solutionbox}


%% ============================================================
\Needspace{5\baselineskip}
\begin{problembox}[Q3-02\quad 损失函数手动计算：回归 vs 分类]
    \textbf{题目描述：}\\
    文档第 3.1.2 和 3.1.3 节分别介绍了分类和回归的损失函数。把公式变成具体数字，是真正理解损失函数的关键一步。

    \vspace{0.3cm}
    \textbf{场景一：回归任务（房价预测）}

    4个样本的真实房价（万元）为 $y = [100, 200, 150, 300]$，模型预测值为 $\hat{y} = [110, 185, 160, 280]$。

    \vspace{0.2cm}
    \textbf{场景二：分类任务（三分类）}

    单个样本，真实类别为第2类（one-hot 标签 $t = [0, 1, 0]$），模型经过 Softmax 后的预测概率为 $\hat{y} = [0.1, 0.7, 0.2]$。

    \vspace{0.3cm}
    \textbf{问题：}
    \begin{enumerate}[(1)]
        \item 对场景一，计算 \textbf{MSE 损失}（结果保留两位小数）：
        \[
        \mathcal{L}_{\text{MSE}} = \frac{1}{N}\sum_{i=1}^{N}(\hat{y}_i - y_i)^2
        \]

        \item 对场景一，计算 \textbf{MAE 损失}（结果保留两位小数）：
        \[
        \mathcal{L}_{\text{MAE}} = \frac{1}{N}\sum_{i=1}^{N}|\hat{y}_i - y_i|
        \]

        \item 对场景二，计算\textbf{交叉熵损失}（结果保留三位小数）：
        \[
        \mathcal{L}_{\text{CE}} = -\sum_{k=1}^{K} t_k \log \hat{y}_k
        \]

        \item 若场景二中预测完全正确（$\hat{y} = [0, 1, 0]$），交叉熵损失理论上应为多少？实际上为什么不能真的输入 $\hat{y}_k = 0$ 或 $\hat{y}_k = 1$（提示：$\log 0 = -\infty$）？PyTorch 的 \texttt{nn.CrossEntropyLoss} 如何解决这个问题？

        \item 比较 (1) 和 (2) 的结果：对于同一组预测误差，MSE 和 MAE 的数值大小有什么规律？MSE 比 MAE 更能"放大"哪种误差？
    \end{enumerate}

    \vspace{0.2cm}
    {\color{gray}\small\textit{来源溯源：[文档第 3 章 3.1.2 分类任务损失函数 / 3.1.3 回归任务损失函数]}}
\end{problembox}

\begin{beginnerbox}[小白必读：交叉熵的直觉——它在衡量什么？]
    \begin{itemize}
        \item \textbf{交叉熵衡量"概率分布的差距"}：真实标签 $t$ 是一个概率分布（one-hot 向量），预测值 $\hat{y}$ 是另一个概率分布（Softmax 输出）。交叉熵越小，两个分布越接近。
        \item \textbf{为什么只有正确类别的项起作用}：one-hot 标签中，错误类别的 $t_k = 0$，对应项 $t_k \log \hat{y}_k = 0$，消失了。所以实际上交叉熵就是 $-\log(\text{正确类别的预测概率})$——预测正确类别的概率越高，Loss 越小，这非常符合直觉。
        \item \textbf{完美预测时 Loss = 0}：若 $\hat{y}_k = 1$（正确类别），$-\log(1) = 0$，Loss 为零——但 Softmax 的输出永远不会真的等于0或1（只能无限接近），所以实践中 Loss 永远大于0。
    \end{itemize}
\end{beginnerbox}

\begin{solutionbox}[参考答案与详细解析]
    \textbf{(1) MSE 损失计算：}

    各样本误差：$10,\ -15,\ 10,\ -20$；误差平方：$100,\ 225,\ 100,\ 400$。
    \[
    \mathcal{L}_{\text{MSE}} = \frac{100 + 225 + 100 + 400}{4} = \frac{825}{4} = \boxed{206.25}\ \text{（万元）}^2
    \]

    \textbf{(2) MAE 损失计算：}

    各样本绝对误差：$10,\ 15,\ 10,\ 20$。
    \[
    \mathcal{L}_{\text{MAE}} = \frac{10 + 15 + 10 + 20}{4} = \frac{55}{4} = \boxed{13.75}\ \text{万元}
    \]

    \textbf{(3) 交叉熵损失计算：}

    真实标签 $t=[0,1,0]$，预测概率 $\hat{y}=[0.1, 0.7, 0.2]$。只有第2类（$k=2$）的 $t_k = 1$ 起作用：
    \[
    \mathcal{L}_{\text{CE}} = -(0 \cdot \log 0.1 + 1 \cdot \log 0.7 + 0 \cdot \log 0.2) = -\log(0.7) \approx \boxed{0.357}
    \]

    \textbf{(4) 数值稳定性处理：}

    理论上完美预测（$\hat{y}_2=1$）时 $\mathcal{L}=-\log(1)=0$。

    但实际中 $\hat{y}_k=0$ 会导致 $\log(0)=-\infty$，产生 \texttt{NaN}，数值计算崩溃。

    \textbf{PyTorch 的解决方案}：\texttt{nn.CrossEntropyLoss} 在内部将 Softmax 和 CrossEntropy 合并，使用 \textbf{log-sum-exp} 技巧在对数域直接计算，同时加入数值稳定性修正（减去最大 logit），避免 $e^{z_k}$ 溢出，也避免 $\log 0$ 问题。这也是为什么该函数要求输入 logits 而非概率。

    \textbf{(5) MSE vs MAE 数值对比与规律：}

    同一组预测误差，MSE = 206.25，MAE = 13.75，\textbf{MSE 远大于 MAE}（单位也不同：MSE 的单位是万元的平方）。

    MSE 特别放大\textbf{大误差}：第4个样本误差为20，在 MAE 中贡献20，在 MSE 中贡献400（是其他误差为10的样本的4倍）。离群点对 MSE 的影响呈平方增长，这既是 MSE 的"缺点"（离群点主导训练），也是优点（对大误差施加更强的惩罚，迫使模型优先改善大误差）。
\end{solutionbox}


%% ============================================================
\Needspace{5\baselineskip}
\begin{problembox}[Q3-03\quad 导数、偏导数与梯度：三个概念辨析]
    \textbf{题目描述：}\\
    文档第 3.2 节介绍了数值微分的基本概念。梯度是深度学习中最核心的数学工具，但"导数""偏导数""梯度"三个词经常被混用。本题把它们彻底分清楚。

    \vspace{0.3cm}
    \textbf{问题：}
    \begin{enumerate}[(1)]
        \item \textbf{概念辨析}：填写下表，区分三个概念：

        \begin{center}
        \begin{tabular}{p{2.2cm} p{4.0cm} p{4.5cm}}
            \toprule
            概念 & 适用对象 & 含义 \\
            \midrule
            导数 $f'(x)$ & 单变量函数 $f(x)$ & \underline{\hspace{4cm}} \\[1.2em]
            偏导数 $\partial f/\partial x_i$ & 多变量函数 $f(x_1,\ldots,x_n)$ & \underline{\hspace{4cm}} \\[1.2em]
            梯度 $\nabla f$ & 多变量函数 $f(x_1,\ldots,x_n)$ & \underline{\hspace{4cm}} \\
            \bottomrule
        \end{tabular}
        \end{center}

        \item \textbf{数值微分计算}：数值微分用差商近似导数。中心差分公式为：
        \[
        f'(x) \approx \frac{f(x+h) - f(x-h)}{2h}
        \]
        对函数 $f(x) = x^3 - 2x$，取 $h = 10^{-4}$，用数值微分估计 $f'(2)$（结果保留四位小数），并与解析导数 $f'(x) = 3x^2 - 2$ 的精确值比较误差。

        \item \textbf{梯度计算}：对二元函数 $f(x_1, x_2) = x_1^2 + 2x_1x_2 + 3x_2^2$，求梯度 $\nabla f$ 的解析表达式，并计算在点 $(1, 2)$ 处的梯度向量。

        \item \textbf{梯度的几何含义}：梯度向量的方向指向函数值增大最快的方向，其大小（模）表示该方向上的变化率。梯度下降沿梯度的\textbf{负方向}更新参数，解释为什么是"负方向"，以及为什么步长（学习率）不能太大或太小。
    \end{enumerate}

    \vspace{0.2cm}
    {\color{gray}\small\textit{来源溯源：[文档第 3 章 3.2 数值微分]}}
\end{problembox}

\begin{beginnerbox}[小白必读：梯度是什么？用最通俗的语言]
    \begin{itemize}
        \item \textbf{导数}（单变量）：函数在某点的"斜率"——往右走一小步，函数值涨了多少？
        \item \textbf{偏导数}（多变量，固定其他变量）：在多维空间中，只沿着某一个坐标轴方向走一小步，函数值涨了多少？
        \item \textbf{梯度}（多变量，所有偏导数打包）：把每个方向的偏导数打包成一个向量 $[\partial f/\partial x_1, \partial f/\partial x_2, \ldots]$。这个向量指向"函数值上升最快的方向"。
        \item \textbf{直觉类比}：你站在一座山上，梯度就是"最陡上坡方向"——梯度下降就是沿着最陡下坡方向走，找山谷（最低点/最小 Loss）。
    \end{itemize}
\end{beginnerbox}

\begin{metaphorbox}[生活比喻：梯度就是地形图上的等高线垂线]
    想象一张地形图，等高线越密说明坡度越陡。

    \begin{itemize}
        \item \textbf{偏导数 $\partial f/\partial x_1$}：只沿东西方向走，看海拔怎么变——等高线在东西方向的变化率；
        \item \textbf{偏导数 $\partial f/\partial x_2$}：只沿南北方向走，看海拔怎么变；
        \item \textbf{梯度向量}：把东西变化率和南北变化率合成一个箭头，这个箭头精确指向"最陡上坡方向"；
        \item \textbf{梯度下降}：沿着这个箭头的反方向走，朝山谷滑去。步子太大（学习率过大）会越过山谷、反弹；步子太小（学习率过小）爬到山谷要走很久。
    \end{itemize}
\end{metaphorbox}

\begin{solutionbox}[参考答案与详细解析]
    \textbf{(1) 概念辨析：}

    \begin{center}
    \begin{tabular}{p{2.2cm} p{4.0cm} p{4.3cm}}
        \toprule
        概念 & 适用对象 & 含义 \\
        \midrule
        导数 $f'(x)$ & 单变量函数 & $x$ 变化一点时，$f(x)$ 变化多快；即函数在该点的切线斜率 \\[0.8em]
        偏导数 $\partial f/\partial x_i$ & 多变量函数 & 固定其他变量，只让 $x_i$ 变化时，$f$ 的变化率 \\[0.8em]
        梯度 $\nabla f$ & 多变量函数 & 所有偏导数组成的向量，指向函数值增大最快的方向 \\
        \bottomrule
    \end{tabular}
    \end{center}

    \textbf{(2) 数值微分计算（$f(x)=x^3-2x$，$x=2$）：}

    \[
    f(2+h) = (2.0001)^3 - 2(2.0001) \approx 8.000600 - 4.0002 = 4.0004
    \]
    \[
    f(2-h) = (1.9999)^3 - 2(1.9999) \approx 7.999400 - 3.9998 = 3.9996
    \]
    \[
    f'(2) \approx \frac{4.0004 - 3.9996}{2 \times 10^{-4}} = \frac{0.0008}{0.0002} = \boxed{10.0000}
    \]
    解析导数：$f'(x) = 3x^2 - 2$，$f'(2) = 12 - 2 = 10$。数值微分结果与解析值完全一致（中心差分的误差为 $O(h^2)$，$h=10^{-4}$ 时误差约 $10^{-8}$，在四位小数内无法体现）。

    \textbf{(3) 梯度计算（$f(x_1,x_2) = x_1^2 + 2x_1x_2 + 3x_2^2$）：}

    \[
    \frac{\partial f}{\partial x_1} = 2x_1 + 2x_2, \quad \frac{\partial f}{\partial x_2} = 2x_1 + 6x_2
    \]
    \[
    \nabla f = \begin{pmatrix} 2x_1+2x_2 \\ 2x_1+6x_2 \end{pmatrix}
    \]
    在点 $(1, 2)$ 处：
    \[
    \nabla f(1,2) = \begin{pmatrix} 2\times1+2\times2 \\ 2\times1+6\times2 \end{pmatrix} = \boxed{\begin{pmatrix} 6 \\ 14 \end{pmatrix}}
    \]

    \textbf{(4) 梯度负方向与学习率的原因：}

    梯度指向 $f$ \textbf{增大}最快的方向。我们的目标是让 Loss \textbf{减小}，所以必须朝梯度的\textbf{负方向}更新参数：$w \leftarrow w - \eta \nabla f$。

    \begin{itemize}
        \item \textbf{学习率 $\eta$ 过大}：每步迈得太大，会跨过 Loss 的最低点，在两侧反复震荡，可能发散（Loss 越来越大）；
        \item \textbf{学习率 $\eta$ 过小}：步子太小，收敛极慢，需要很多轮才能到达最优解，训练时间过长；
        \item \textbf{最优学习率}：使每步更新都朝最优解迈进而不震荡，通常通过实验或学习率衰减策略找到。
    \end{itemize}
\end{solutionbox}


%% ============================================================
\Needspace{5\baselineskip}
\begin{problembox}[Q3-04\quad 神经网络的梯度计算]
    \textbf{题目描述：}\\
    文档第 3.3 节介绍了如何计算神经网络的梯度。对于一个单层神经网络（线性层 + Sigmoid + 二分类交叉熵损失），我们来手动用数值微分估算梯度。

    \vspace{0.3cm}
    网络结构：$z = w \cdot x + b$，$\hat{y} = \sigma(z)$，$L = -[y\log\hat{y} + (1-y)\log(1-\hat{y})]$

    当前参数：$w = 0.5,\ b = 0,\ x = 2.0,\ y = 1$（真实标签），$h = 10^{-4}$。

    \vspace{0.3cm}
    \textbf{问题：}
    \begin{enumerate}[(1)]
        \item \textbf{前向传播基准值}：计算当前参数下的 $z$、$\hat{y}$、$L$（结果保留四位小数）。

        \item \textbf{数值微分估算 $\partial L / \partial w$}：
        \begin{enumerate}[a.]
            \item 将 $w$ 增加 $h$（即 $w_+ = 0.5 + 10^{-4}$），计算 $L(w_+)$；
            \item 将 $w$ 减少 $h$（即 $w_- = 0.5 - 10^{-4}$），计算 $L(w_-)$；
            \item 用中心差分 $\partial L/\partial w \approx (L(w_+)-L(w_-))/(2h)$ 估算梯度（结果保留四位小数）。
        \end{enumerate}

        \item \textbf{数值微分估算 $\partial L / \partial b$}：同理对 $b$ 做数值微分，估算 $\partial L / \partial b$。

        \item \textbf{梯度的物理含义}：根据计算结果，$\partial L / \partial w$ 的正负说明了什么？如果使用学习率 $\eta = 0.1$，执行一步梯度下降后，新的 $w$ 和 $b$ 分别是多少？

        \item \textbf{解析梯度验证}：对于 Sigmoid + 二分类交叉熵的组合，可以推导出：
        \[
        \frac{\partial L}{\partial w} = (\hat{y} - y) \cdot x, \quad \frac{\partial L}{\partial b} = \hat{y} - y
        \]
        用上述公式计算解析梯度，与数值微分结果对比，误差是多少？
    \end{enumerate}

    \vspace{0.2cm}
    {\color{gray}\small\textit{来源溯源：[文档第 3 章 3.3 神经网络的梯度计算]}}
\end{problembox}

\begin{beginnerbox}[小白必读：为什么要"手算"梯度这件事？]
    \begin{itemize}
        \item \textbf{建立直觉}：手算一遍梯度，你才能真正理解"梯度告诉你什么"——它的正负说明参数应该增大还是减小，它的大小说明调整幅度应该多大。
        \item \textbf{理解数值微分的用途}：数值微分（有限差分）是验证解析梯度是否正确的"金标准"——虽然计算量大，但不需要推导公式，只要能前向计算损失就行。
        \item \textbf{感受链式法则的威力}：从本题可以看出，即使是最简单的单节点网络，数值微分也需要多次前向传播。真实网络有数百万参数，解析梯度（反向传播）一次性算出所有梯度，效率高出数量级。
    \end{itemize}
\end{beginnerbox}

\begin{solutionbox}[参考答案与详细解析]
    \textbf{(1) 前向传播基准值（$w=0.5,\ b=0,\ x=2,\ y=1$）：}
    \[
    z = 0.5 \times 2 + 0 = 1.0
    \]
    \[
    \hat{y} = \sigma(1.0) = \frac{1}{1+e^{-1}} = \frac{1}{1+0.3679} \approx 0.7311
    \]
    \[
    L = -\log(0.7311) \approx \boxed{0.3133}
    \]

    \textbf{(2) 数值微分估算 $\partial L/\partial w$：}

    \textbf{(a)} $w_+ = 0.5001$：$z_+ = 0.5001 \times 2 = 1.0002$，$\hat{y}_+ = \sigma(1.0002) \approx 0.7312$，$L_+ = -\log(0.7312) \approx 0.3132$

    \textbf{(b)} $w_- = 0.4999$：$z_- = 0.4999 \times 2 = 0.9998$，$\hat{y}_- = \sigma(0.9998) \approx 0.7310$，$L_- = -\log(0.7310) \approx 0.3134$

    \textbf{(c)} 中心差分：
    \[
    \frac{\partial L}{\partial w} \approx \frac{0.3132 - 0.3134}{2 \times 10^{-4}} = \frac{-0.0002}{0.0002} = \boxed{-0.5378}
    \]
    （精确数值视舍入而定，约为 $-0.5378$）

    \textbf{(3) 数值微分估算 $\partial L/\partial b$：}

    $b_\pm = \pm 10^{-4}$，$z_\pm = 1.0 \pm 10^{-4}$，差值引起的 $L$ 变化为 $\hat{y}-1$ 的量级：
    \[
    \frac{\partial L}{\partial b} \approx \boxed{-0.2689}
    \]

    \textbf{(4) 梯度物理含义与一步更新：}

    $\partial L/\partial w \approx -0.5378 < 0$：表示增大 $w$ 会使 Loss 减小，即 $w$ 目前偏小，应增大。梯度下降沿负梯度方向（即正方向）更新 $w$。

    一步梯度下降（$\eta=0.1$）：
    \[
    w_{\text{new}} = 0.5 - 0.1 \times (-0.5378) = 0.5 + 0.0538 = \boxed{0.5538}
    \]
    \[
    b_{\text{new}} = 0 - 0.1 \times (-0.2689) = \boxed{0.0269}
    \]

    \textbf{(5) 解析梯度验证：}

    \[
    \frac{\partial L}{\partial w} = (\hat{y} - y) \cdot x = (0.7311 - 1) \times 2 = -0.2689 \times 2 = \boxed{-0.5378}
    \]
    \[
    \frac{\partial L}{\partial b} = \hat{y} - y = 0.7311 - 1 = \boxed{-0.2689}
    \]

    解析梯度与数值微分结果完全吻合（误差在 $10^{-8}$ 量级，远小于四位小数精度），验证正确。\checkmark

    \begin{keybox}[关键结论：Sigmoid + 二元交叉熵的梯度极简公式]
        \[
        \frac{\partial L}{\partial w} = (\hat{y} - y) \cdot x, \quad \frac{\partial L}{\partial b} = \hat{y} - y
        \]
        含义：预测概率 $\hat{y}$ 和真实标签 $y$ 之差就是误差信号，直接乘以对应输入 $x$ 即为权重梯度。预测越准（$\hat{y}\to y$），梯度越小，更新越慢——这正是我们希望的行为。
    \end{keybox}
\end{solutionbox}


%% ============================================================
\Needspace{5\baselineskip}
\begin{problembox}[Q3-05\quad 梯度下降三兄弟：BGD、SGD、Mini-batch GD]
    \textbf{题目描述：}\\
    文档第 3.4 节介绍了随机梯度下降法。"梯度下降"并不是一种固定的算法，根据每次更新时使用的数据量，分为三种变体，各有优缺点。

    \vspace{0.3cm}
    \textbf{问题：}
    \begin{enumerate}[(1)]
        \item 填写下表，对比三种梯度下降变体：

        \begin{center}
        \begin{tabular}{p{2.5cm} p{3.0cm} p{3.0cm} p{3.2cm}}
            \toprule
            方法 & 每次更新使用的数据量 & 主要优点 & 主要缺点 \\
            \midrule
            批量梯度下降（BGD） & 全部训练集 $N$ 个样本 & \underline{\hspace{2.5cm}} & \underline{\hspace{2.5cm}} \\[1.5em]
            随机梯度下降（SGD） & 单个样本 & \underline{\hspace{2.5cm}} & \underline{\hspace{2.5cm}} \\[1.5em]
            小批量梯度下降（Mini-batch） & $B$ 个样本（$1 < B < N$） & \underline{\hspace{2.5cm}} & \underline{\hspace{2.5cm}} \\
            \bottomrule
        \end{tabular}
        \end{center}

        \item \textbf{Epoch 与 Iteration 的区别}：设训练集有 $N=10000$ 个样本，batch size $B=100$，训练 20 个 epoch。计算：
        \begin{enumerate}[a.]
            \item 每个 epoch 包含多少次参数更新（iterations）？
            \item 总共进行多少次参数更新？
            \item 总共处理了多少个样本次数（样本$\times$次数）？
        \end{enumerate}

        \item \textbf{为什么 SGD 有时比 BGD 收敛更快}：从"逃离局部极小值"和"更新频率"两个角度解释。

        \item \textbf{Batch Size 的选择}：以下三种说法中，哪些是正确的？请逐一判断并说明理由：
        \begin{enumerate}[a.]
            \item "Batch size 越大，模型泛化性能越好"
            \item "Batch size 过小，GPU 利用率低，训练慢"
            \item "Batch size 的选择和学习率应该一起调整"
        \end{enumerate}

        \item \textbf{数据打乱（Shuffle）}：为什么每个 epoch 开始前通常要对训练数据进行 shuffle（随机打乱）？如果不打乱，会出现什么问题？
    \end{enumerate}

    \vspace{0.2cm}
    {\color{gray}\small\textit{来源溯源：[文档第 3 章 3.4 随机梯度下降法（SGD）]}}
\end{problembox}

\begin{metaphorbox}[生活比喻：三种梯度下降像三种市场调研方式]
    你是一家公司的产品经理，想知道产品需要改进哪里（梯度方向）：

    \begin{itemize}
        \item \textbf{BGD（全量调研）}：把所有100万用户都调查一遍，再做决策。结果最准确，但等完所有反馈才能改一次，太慢了；
        \item \textbf{SGD（随机抽一人）}：随机问一个用户的意见，立刻就改。改得飞快，但这个用户可能是异类，反馈噪声很大，产品走走停停；
        \item \textbf{Mini-batch（抽一小组）}：每次随机抽100人调研，综合他们的意见再改。既比全量快得多，反馈又比单人靠谱——实践中几乎所有深度学习都用这种方式。
    \end{itemize}
\end{metaphorbox}

\begin{solutionbox}[参考答案与详细解析]
    \textbf{(1) 三种梯度下降对比表：}

    \begin{center}
    \begin{tabular}{p{2.5cm} p{3.0cm} p{3.0cm} p{2.8cm}}
        \toprule
        方法 & 数据量 & 主要优点 & 主要缺点 \\
        \midrule
        BGD & 全部 $N$ 个 & 梯度估计准确，收敛稳定 & 内存需求大，每步耗时长 \\[0.8em]
        SGD & 1个 & 更新频率极高，可逃离局部极小 & 梯度噪声大，Loss震荡，难并行 \\[0.8em]
        Mini-batch & $B$ 个 & 平衡速度与准确性，可充分利用GPU并行 & 需调 batch size 超参数 \\
        \bottomrule
    \end{tabular}
    \end{center}

    \textbf{(2) Epoch 与 Iteration 计算：}

    \begin{enumerate}[a.]
        \item 每个 epoch 的 iterations = $N / B = 10000 / 100 = \boxed{100}$ 次
        \item 总参数更新次数 = $100 \times 20 = \boxed{2000}$ 次
        \item 总处理样本次数 = $10000 \times 20 = \boxed{200000}$ 个样本次数
    \end{enumerate}

    \textbf{(3) SGD 有时比 BGD 收敛更快的原因：}

    \begin{itemize}
        \item \textbf{逃离局部极小值}：BGD 每次沿精确梯度更新，容易陷在局部极小值（停止更新）。SGD 的梯度噪声相当于给参数更新加了"随机扰动"，有机会跳出局部极小值，找到更好的解；
        \item \textbf{更新频率更高}：BGD 处理完全部 $N$ 个样本才更新一次，SGD 处理1个样本就更新一次。在相同的数据处理量下，SGD 更新了 $N$ 次，参数已经朝最优解迈进了 $N$ 步，而 BGD 只走了1步。
    \end{itemize}

    \textbf{(4) Batch Size 三种说法的判断：}

    \begin{enumerate}[a.]
        \item \textbf{错误}。研究表明 batch size 越大，模型往往泛化性能\textbf{越差}（倾向于收敛到"尖锐的"极小值，泛化能力弱）。小 batch size 引入的噪声有助于找到"平坦的"极小值，泛化更好；
        \item \textbf{正确}。GPU 的并行计算需要足够的数据量才能充分利用显存带宽。batch size = 1 时，GPU 大部分时间在等待数据，利用率极低；
        \item \textbf{正确}。增大 batch size 相当于减小了梯度的方差，等效于"学习率变小"。实践中有一个经验规律：\textbf{线性缩放规则}——batch size 增大 $k$ 倍时，学习率也应增大 $k$ 倍，以保持等效的参数更新幅度。
    \end{enumerate}

    \textbf{(5) Shuffle 的必要性：}

    如果数据按类别排列（如前1000张猫、后1000张狗），不打乱时每个 mini-batch 只包含单一类别，梯度估计严重偏置——模型只对当前批次的类别优化，损害了泛化性。此外，若数据有时间或顺序上的规律，不打乱会让模型学到顺序特征而非真实规律。

    每个 epoch 前 shuffle 能保证每个 mini-batch 是训练集的随机子集，梯度估计更接近全数据集的真实梯度，训练更稳定，泛化更好。
\end{solutionbox}


%% ============================================================
\Needspace{5\baselineskip}
\begin{problembox}[Q3-06\quad 综合代码实现：完整训练循环]
    \textbf{题目描述：}\\
    文档第 3.5 节给出了神经网络学习的综合代码实现。以下是一个用 NumPy 实现的完整单层神经网络训练代码骨架（二分类任务），请阅读并补全关键步骤。

    \begin{lstlisting}[language=Python]
import numpy as np

# ===== 辅助函数 =====
def sigmoid(z):
    return 1 / (1 + np.exp(-z))

def binary_cross_entropy(y_hat, y):
    eps = 1e-8  # 防止 log(0)
    return -np.mean(y * np.log(y_hat + eps)
                    + (1 - y) * np.log(1 - y_hat + eps))

def numerical_gradient(loss_fn, param, h=1e-4):
    grad = np.zeros_like(param)
    for i in range(param.size):
        param.flat[i] += h
        lp = loss_fn()
        param.flat[i] -= 2 * h
        lm = loss_fn()
        grad.flat[i] = (lp - lm) / (2 * h)
        param.flat[i] += h
    return grad

# ===== 数据与初始化 =====
np.random.seed(0)
X = np.random.randn(100, 3)    # 100 个样本，3 个特征
y = (X[:, 0] + X[:, 1] > 0).astype(float)  # 真实标签

W = np.random.randn(3, 1) * 0.01  # 权重
b = np.zeros(1)                    # 偏置
lr = 0.1

# ===== 训练循环 =====
for epoch in range(50):
    # ------ 步骤A：前向传播 ------
    Z     = _______________          # 线性变换：[100,3] @ [3,1] + [1]
    Y_hat = _______________          # Sigmoid 激活

    # ------ 步骤B：计算损失 ------
    loss = _______________

    # ------ 步骤C：计算梯度（用解析梯度）------
    dZ = _______________             # Sigmoid+CE 组合梯度：Y_hat - y（注意形状）
    dW = _______________             # dL/dW
    db = _______________             # dL/db

    # ------ 步骤D：参数更新 ------
    W -= _______________
    b -= _______________

    if epoch % 10 == 0:
        print(f"Epoch {epoch:3d} | Loss: {loss:.4f}")
    \end{lstlisting}

    \textbf{问题：}
    \begin{enumerate}[(1)]
        \item 补全步骤A、B、C、D中的所有空白（共7处），写出完整的填写内容，并注意矩阵形状的正确性。

        \item \textbf{形状追踪}：填写以下各变量的形状：
        \begin{enumerate}[a.]
            \item \texttt{X}：\underline{\hspace{3cm}}，\texttt{W}：\underline{\hspace{3cm}}，\texttt{b}：\underline{\hspace{3cm}}
            \item \texttt{Z}（线性变换输出）：\underline{\hspace{3cm}}
            \item \texttt{Y\_hat}（Sigmoid 输出）：\underline{\hspace{3cm}}
            \item \texttt{dZ}（梯度）：\underline{\hspace{3cm}}
            \item \texttt{dW}（权重梯度）：\underline{\hspace{3cm}}，\texttt{db}（偏置梯度）：\underline{\hspace{3cm}}
        \end{enumerate}

        \item \textbf{步骤C 的推导依据}：解析梯度 $\mathrm{d}W = X^T \cdot \mathrm{d}Z / N$ 来自哪个推导结果（参考 Q3-04 和 Q4-05）？为什么要除以 $N$？

        \item \textbf{训练过程观察}：如果每10个 epoch 打印一次 Loss，你期望 Loss 的变化趋势是什么？如果 Loss 不降反升，最可能的原因是什么？

        \item \textbf{改进方向}：这段代码使用了数值微分还是解析梯度（步骤C）？如果改用 \texttt{numerical\_gradient} 函数替代步骤C，训练时间会有什么变化？（提示：考虑参数总数和每次数值微分需要的前向传播次数）
    \end{enumerate}

    \vspace{0.2cm}
    {\color{gray}\small\textit{来源溯源：[文档第 3 章 3.5 综合代码实现]}}
\end{problembox}

\begin{beginnerbox}[小白必读：训练循环的四步节奏，永远不会变]
    \begin{itemize}
        \item \textbf{第一步，前向传播}：数据流过网络，计算预测值；
        \item \textbf{第二步，计算 Loss}：预测值和真实标签比较，得到一个数字；
        \item \textbf{第三步，计算梯度}：Loss 对每个参数的偏导数（方向盘告诉你往哪转）；
        \item \textbf{第四步，更新参数}：沿梯度负方向走一步（$w \leftarrow w - \eta \cdot \nabla_w L$）。
    \end{itemize}
    这四步就是深度学习训练的最小闭环。无论 PyTorch 的训练循环写得多复杂，拆开来看都是这四步：\texttt{forward} $\to$ \texttt{loss} $\to$ \texttt{backward} $\to$ \texttt{optimizer.step()}。
\end{beginnerbox}

\begin{solutionbox}[参考答案与详细解析]
    \textbf{(1) 完整填写内容：}

    \begin{lstlisting}[language=Python]
# 步骤A：前向传播
Z     = X @ W + b                          # [100,3]@[3,1]+[1] = [100,1]
Y_hat = sigmoid(Z)                         # [100,1]

# 步骤B：计算损失
loss  = binary_cross_entropy(Y_hat, y.reshape(-1, 1))

# 步骤C：计算解析梯度
N  = X.shape[0]                            # 样本数 100
dZ = (Y_hat - y.reshape(-1, 1)) / N       # [100,1]，除以N取均值
dW = X.T @ dZ                             # [3,100]@[100,1] = [3,1]
db = np.sum(dZ, axis=0)                   # [1]，对样本求和

# 步骤D：参数更新
W -= lr * dW
b -= lr * db
    \end{lstlisting}

    \textbf{(2) 形状追踪：}

    \begin{center}
    \begin{tabular}{lll}
        \toprule
        变量 & 形状 & 说明 \\
        \midrule
        \texttt{X} & \texttt{[100, 3]} & 100个样本，3个特征 \\
        \texttt{W} & \texttt{[3, 1]} & 3个输入特征，1个输出 \\
        \texttt{b} & \texttt{[1]} & 标量偏置 \\
        \texttt{Z} & \texttt{[100, 1]} & 每个样本的线性输出 \\
        \texttt{Y\_hat} & \texttt{[100, 1]} & Sigmoid 后的预测概率 \\
        \texttt{dZ} & \texttt{[100, 1]} & 每个样本的误差信号（除以N） \\
        \texttt{dW} & \texttt{[3, 1]} & 与 $W$ 形状相同 \\
        \texttt{db} & \texttt{[1]} & 与 $b$ 形状相同 \\
        \bottomrule
    \end{tabular}
    \end{center}

    \textbf{(3) $\mathrm{d}W = X^T \cdot \mathrm{d}Z / N$ 的推导依据：}

    来自 Q3-04 中的单样本结论 $\partial L/\partial w = (\hat{y}-y)\cdot x$ 推广到批量情形，以及 Q4-05 中的 Affine 层矩阵梯度推导 $\partial L/\partial W = X^T \delta_Y$。

    除以 $N$ 是因为 \texttt{binary\_cross\_entropy} 对 $N$ 个样本取了均值（\texttt{np.mean}），对应地，梯度也要对 $N$ 个样本取均值——即 $\mathrm{d}Z$ 中已除以 $N$，$\mathrm{d}W = X^T \cdot \mathrm{d}Z$ 自然也是均值梯度。这保证了损失尺度和梯度尺度在不同 batch size 下保持一致。

    \textbf{(4) 训练过程观察：}

    正常情况下，Loss 应\textbf{单调下降}（可能有小幅波动），最终趋于平稳。

    Loss 不降反升的最常见原因：
    \begin{itemize}
        \item \textbf{学习率过大}：每步更新幅度太大，越过最优点，Loss 震荡或发散——最常见原因，首先尝试将 $\eta$ 减小10倍；
        \item \textbf{梯度计算错误}：符号写反（用了 $+\eta$ 而不是 $-\eta$），或梯度公式有误；
        \item \textbf{数值溢出}：$e^z$ 在 $z$ 很大时溢出为 \texttt{inf}，导致 Loss 变成 \texttt{NaN}。
    \end{itemize}

    \textbf{(5) 解析梯度 vs 数值微分的效率：}

    步骤C 使用的是\textbf{解析梯度}（直接用公式计算，一次前向传播后计算即完成）。

    若改用 \texttt{numerical\_gradient}：本例共有 $3\times1 + 1 = 4$ 个参数，每个参数需要2次前向传播，共需 $4 \times 2 = 8$ 次前向传播才能算出一次梯度。对真实网络（如参数量 $10^7$），数值微分需要 $2\times10^7$ 次前向传播，而解析梯度（反向传播）只需1次，\textbf{效率差距约 $2\times10^7$ 倍}——数值微分在实际训练中完全不可行。

    \begin{keybox}[完整训练循环模板（可直接复用）]
        \begin{lstlisting}[language=Python]
for epoch in range(num_epochs):
    # 1. 前向传播
    Z     = X @ W + b
    Y_hat = sigmoid(Z)

    # 2. 计算损失
    loss = binary_cross_entropy(Y_hat, y)

    # 3. 计算梯度（解析梯度）
    dZ = (Y_hat - y) / N
    dW = X.T @ dZ
    db = np.sum(dZ, axis=0)

    # 4. 更新参数
    W -= lr * dW
    b -= lr * db
        \end{lstlisting}
    \end{keybox}
\end{solutionbox}

\begin{appbox}[现实应用场景]
    \textbf{Mini-batch SGD——支撑起整个深度学习时代的训练策略：}\\
    本章看似"简单"的随机梯度下降，是所有现代大模型训练的基础。GPT 系列在预训练阶段，使用 batch size 数百万 token、学习率精心调节的 AdamW（SGD 的改进版），在数千个 GPU 上同时运行 mini-batch 梯度下降。理解"每个 batch 算一次梯度、更新一次参数"这个核心循环，就理解了深度学习训练的本质——无论规模大到何种程度，这个闭环从未改变。
\end{appbox}

% ===== 本章小结 =====
\section{本章小结：学习的三件事全打通了}

学完这 6 道题，整章内容可以浓缩成这一句话：

\begin{center}
    \textbf{用损失函数量化错误 $\xrightarrow{\text{梯度}}$ 知道往哪改 $\xrightarrow{\text{SGD}}$ 一步步变好}
\end{center}

\begin{itemize}
    \item \textbf{损失函数}：回归用 MSE/MAE/Huber，分类用交叉熵；选错方向，优化就跑偏；
    \item \textbf{梯度}：导数的多变量推广，方向指向 Loss 增大最快处；沿负梯度更新参数，Loss 就会减小；
    \item \textbf{数值微分}：用差商近似导数，是验证梯度的金标准，但计算量与参数数量成正比，实际训练不可用；
    \item \textbf{SGD 三兄弟}：BGD 准但慢，SGD 快但噪声大，Mini-batch 折中两者，是深度学习的实际标准；
    \item \textbf{训练循环四步}：前向 $\to$ 计算Loss $\to$ 计算梯度 $\to$ 更新参数，这个闭环无论网络多深、框架多复杂，永远不变。
\end{itemize}

\begin{metaphorbox}[给零基础读者的一句总比喻]
    \textbf{如果把训练神经网络比作学打靶：}
    \begin{itemize}
        \item \textbf{损失函数} = 报靶员，告诉你偏了几环（MSE 对大偏差惩罚重，MAE 一视同仁，交叉熵对"完全没射中"惩罚最重）；
        \item \textbf{梯度} = 教练，告诉你"弓往左转、往右转还是往上抬"——梯度的方向就是需要调整的方向（取反）；
        \item \textbf{学习率} = 每次调整的幅度：调太猛（大学习率），可能从左偏变成右偏；调太保守（小学习率），射100箭还没对准；
        \item \textbf{Mini-batch SGD} = 不等你把靶场所有人的意见都收集完再调，每次随机问一小组人，马上调，虽然每次不完美，但总体射得越来越准——快！
    \end{itemize}
\end{metaphorbox}

\end{document}
